\chapter{最小 CPU 与控制 trace}

前面已经具备构造处理器的必要部件：组合逻辑、状态元件、时钟边界、ALU、数据通路和控制器。本章只做一件事：把这些部件合成为一台能够取指、译码、执行、写回和更新 PC 的最小 CPU。主存、总线、设备和经典芯片会在后续章节展开，本章先把 CPU 内部的状态边界和控制 trace 写清楚。

CPU 可以视为一组同步硬件闭环：PC 指出下一条编码地址，取指路径取回指令 bit，译码逻辑生成控制信号，数据通路读取和计算，状态元件在时钟边界提交结果。只要能为一条指令写出 PC、IR、寄存器堆、ALU、写回 mux、标志位和控制器的完整 trace，就已经能判断这台最小 CPU 是否真的会执行程序。

\begin{keyidea}
最小 CPU 不是“ALU 加寄存器”的粗线图，而是一组按周期提交的状态转换。每条指令都必须说明：本周期读哪些旧状态，组合路径经过哪些部件，哪些控制线有效，时钟沿会提交哪些新状态，异常或等待是否会阻止提交。
\end{keyidea}

\begin{pdfdiagram}
  \node[hw state, minimum width=1.7cm] (pc) at (0,0.75) {PC};
  \node[hw compute, minimum width=2.0cm] (imem) at (2.55,0.75) {指令存储};
  \node[hw state, minimum width=1.55cm] (ir) at (5.05,0.75) {IR};
  \node[hw control, minimum width=1.9cm] (ctrl) at (7.55,0.75) {控制器};
  \node[hw state, minimum width=2.0cm] (rf) at (2.55,-1.0) {寄存器堆};
  \node[hw compute, minimum width=1.7cm] (alu) at (5.05,-1.0) {ALU};
  \node[hw compute, minimum width=2.0cm] (wb) at (7.55,-1.0) {写回 mux};
  \draw[hw bus] (pc) -- node[above=2pt,hw label,fill=white,inner sep=1pt]{取指地址} (imem);
  \draw[hw bus] (imem) -- node[above=2pt,hw label,fill=white,inner sep=1pt]{编码 bit} (ir);
  \draw[hw bus] (ir) -- node[above=2pt,hw label,fill=white,inner sep=1pt]{字段} (ctrl);
  \draw[hw return] (ctrl.south) |- node[pos=0.25,right,hw label,fill=white,inner sep=1pt]{控制线} (5.05,-0.18) -| (alu.north);
  \draw[hw bus] (rf) -- node[above=2pt,hw label,fill=white,inner sep=1pt]{源操作数} (alu);
  \draw[hw bus] (alu) -- node[above=2pt,hw label,fill=white,inner sep=1pt]{结果/标志} (wb);
  \draw[hw return] (wb.east) -- ++(0.55,0) |- (rf.south) node[pos=0.25,below=2pt,hw label,fill=white,inner sep=1pt]{写回};
  \draw[hw return] (ctrl.north) |- (3.85,1.75) -| (pc.north) node[pos=0.2,above=2pt,hw label,fill=white,inner sep=1pt]{下一 PC};
  \node[hw label, text width=9.4cm] at (4.1,-2.05) {最小 CPU 的核心验收对象是 trace：PC 取指，IR 保存编码，控制器打开数据通路，寄存器和 PC 在时钟沿提交。};
\end{pdfdiagram}
\diagramnote{最小 CPU 总图。箭头表示一次指令从取指到写回的证据链，反馈线表示 PC 和寄存器堆在时钟边界更新。}

\section{一、从动作编码到最小 CPU}

若希望同一套数据通路在不同时间做不同动作，就需要把“这一次做什么”也表示成 bit。动作编码可以先写成最小形式：

\begin{lstlisting}
opcode | dst | src_a | src_b
\end{lstlisting}

\code{opcode} 表示动作类型，例如 ADD、SUB、AND；\code{dst} 表示写入哪个寄存器；\code{src\_a} 和 \code{src\_b} 表示读取哪两个寄存器。控制器读取这些字段后，把它们译成 ALU op、读地址、写地址、写使能和写回选择。编码本身仍然是硬件输入，不是文字命令。

若有一组动作编码连续放在存储结构中，硬件还需要知道当前取哪一条。这个位置由 PC 保存。PC 本质上也是寄存器，只是它保存的是“下一条动作编码的地址”。

\begin{longtable}{L{0.24\linewidth}L{0.52\linewidth}}
\toprule
部件 & 作用 \\
\midrule
PC & 保存下一条编码地址 \\
指令存储器 & 根据 PC 输出当前编码 \\
控制器 & 根据编码字段产生控制信号 \\
寄存器阵列 & 保存通用数据 \\
ALU & 做整数和逻辑运算 \\
选择器网络 & 选择 ALU 输入、写回来源和下一 PC \\
\bottomrule
\end{longtable}

顺序执行时，下一 PC 可以由加法器产生：

\begin{lstlisting}
pc_next = pc_current + instruction_size
\end{lstlisting}

如果每条编码占 4 个字节，\code{instruction\_size} 就是 4；如果是可变长度指令，下一 PC 的计算会更复杂。关键事实是：顺序执行不是抽象的“下一行”，而是 PC 这个状态寄存器保存了新的地址。

最小单周期模型可以这样走：

\begin{longtable}{L{0.18\linewidth}L{0.52\linewidth}L{0.18\linewidth}}
\toprule
阶段 & 硬件动作 & 边界 \\
\midrule
取指 & PC 送入指令存储器，得到编码 & 组合路径 \\
译码 & 控制器产生 mux、ALU、写使能信号 & 组合路径 \\
执行 & 寄存器阵列读值，ALU 计算 & 组合路径 \\
写回 & 结果在时钟沿写入寄存器或 PC & 状态边界 \\
\bottomrule
\end{longtable}

用一条普通运算编码走读：当前位置 PC=100，指令存储器在地址 100 处保存的是：

\begin{lstlisting}
ADD dst=R3, src_a=R1, src_b=R2
\end{lstlisting}

周期开始时，PC 输出稳定为 100。这个地址进入指令存储器，存储器经过地址译码后把这条编码输出。控制器看到 opcode=ADD，于是把 \code{alu\_op} 设为 ADD，把寄存器阵列的两个读地址分别设为 R1 和 R2，把写地址设为 R3，把 \code{wb\_sel} 设为 ALU，把 \code{reg\_write} 设为 1，同时让下一 PC 选择 PC 加指令长度。

这些控制信号稳定后，寄存器阵列的两个读端口输出 R1 和 R2 当前的值。它们经过 ALU 输入选择器进入 ALU，加法路径产生结果。结果再经过写回选择器，到达寄存器阵列的写数据端口。另一边，PC 加法器准备下一条指令地址，PC 选择器把它送到 PC 的下一值输入。只有到时钟沿，两个状态变化才同时提交：R3 保存 ALU 结果，PC 保存下一条地址。

分支也可以归结为下一 PC 选择。顺序情况下，下一 PC 是 PC 加指令长度；条件分支则由 ALU 比较结果和选择器决定下一 PC 来源：

\begin{lstlisting}
branch_taken = condition_met AND branch_op
next_pc = mux(branch_taken, pc_sequential, branch_target)
\end{lstlisting}

分支的硬件本质是改变 PC 这个寄存器的下一值。下一周期取到哪条编码，完全由这个新 PC 决定。

\topic{把指令周期写成可验收账本}经典教材讲 CPU 时，常把一条指令拆成取指、译码、执行、访存、写回和下一 PC 更新。这个拆分不是为了背阶段名称，而是为了把每个阶段对应到真实硬件路径。最小单周期 CPU 可以在一个长周期内完成所有路径；多周期 CPU 可以把这些路径分到多个时钟；流水线 CPU 可以让多条指令在不同阶段重叠推进。三者的共同点是：每个阶段都要说明源状态、组合路径、目标状态和异常边界。

\begin{longtable}{L{0.14\linewidth}L{0.34\linewidth}L{0.28\linewidth}L{0.14\linewidth}}
\toprule
阶段 & 主要硬件动作 & 关键控制信号 & 提交边界 \\
\midrule
IF 取指 & PC 送到指令存储器或 I-cache，取回指令 bit & \code{pc_out}、\code{imem_read}、\code{ir_write} & IR 或取指缓冲 \\
ID 译码 & opcode、寄存器号、立即数字段被解释成控制线 & \code{decode}、\code{imm_kind}、\code{read_addr} & 可组合，也可锁存 \\
EX 执行 & ALU 做加减、逻辑、比较或地址形成 & \code{alu_op}、\code{a_sel}、\code{b_sel} & 输出和标志 \\
MEM 访存 & 对数据存储器、cache 或 MMIO 发起读写事务 & \code{mem_read}、\code{mem_write}、\code{byte_en} & MDR 或外部响应 \\
WB 写回 & ALU 结果或读回数据写入寄存器堆 & \code{wb_sel}、\code{reg_write} & 目标寄存器 \\
PC 更新 & 选择顺序 PC、分支目标、跳转目标或异常入口 & \code{pc_sel}、\code{pc_write} & PC \\
\bottomrule
\end{longtable}

这张表可以直接变成验收清单。若某条 LOAD 指令读到了错误数据，不应只问“内存为什么错”，而要逐项检查：PC 是否取到正确编码，译码是否识别为 LOAD，ALU 是否形成正确地址，MEM 阶段是否访问正确目标，读回数据是否进入 MDR，WB 阶段是否选择了 MDR 而不是 ALUOut。这样的排查方式比凭感觉改控制器可靠得多。

\topic{控制字把“会执行指令”落到可接线信号}教学 CPU 最容易写成口号：取指、译码、执行、写回。真正能搭出来的版本必须有控制字。控制字不是高级软件结构，而是一组在某个周期稳定的控制线集合。它决定寄存器是否写入、ALU 选哪个操作、存储器是否读写、PC 从哪里取下一值。若控制字没有定义清楚，CPU 即使图上连通，也无法稳定执行程序。

\begin{longtable}{L{0.18\linewidth}L{0.2\linewidth}L{0.2\linewidth}L{0.2\linewidth}L{0.12\linewidth}}
\toprule
控制对象 & ADD & LOAD & STORE & JZ \\
\midrule
ALU 操作 & 加法或指定算术 & 基址加偏移 & 基址加偏移 & 比较或测试零 \\
寄存器读 & 两个源寄存器 & 基址寄存器 & 基址和数据寄存器 & 条件源寄存器 \\
存储器读 & 无 & 有 & 无 & 无 \\
存储器写 & 无 & 无 & 有 & 无 \\
写回选择 & ALU 结果 & 存储器数据 & 无 & 无 \\
PC 选择 & 顺序 PC & 顺序 PC & 顺序 PC & 条件目标或顺序 PC \\
\bottomrule
\end{longtable}

若把这些控制对象压成 bit，可以得到类似下面的教学控制字。具体 bit 位不重要，重要的是每一位都能对应到实际连线或选择器输入。

\begin{lstlisting}
control word fields:
  pc_sel      // sequential, branch_target, jump_target, exception
  reg_write   // enable register-file write port
  wb_sel      // ALU, memory, pc_plus_size
  alu_op      // add, sub, and, or, compare
  alu_b_sel   // register or immediate
  mem_read    // request data read
  mem_write   // request data write
  byte_en     // selected byte lanes
  flags_write // update Zero/Carry/Overflow/Sign
\end{lstlisting}

\topic{硬连线控制与微程序控制是两种组织控制字的方法}控制器可以直接由 opcode、状态位和当前周期译码出控制线，这称为硬连线控制。它速度快，结构适合简单指令集和高频流水线，但修改指令行为需要改组合逻辑。另一种方法是把控制字存在控制存储器中，由微程序计数器逐条取出微指令；每条机器指令对应一段微操作序列。这种方法更像“硬件内部的小程序”，适合组织复杂指令、多周期总线访问和异常入口，但控制存储器和微序列会增加延迟。

\begin{longtable}{L{0.18\linewidth}L{0.34\linewidth}L{0.38\linewidth}}
\toprule
控制方式 & 适合场景 & 验收问题 \\
\midrule
硬连线控制 & 指令格式规整、周期少、目标频率高 & opcode、状态位和时序状态是否唯一决定控制线 \\
微程序控制 & 指令复杂、多周期动作多、需要复用数据通路 & 微地址如何选择、分支微指令如何跳转、异常如何打断 \\
混合方式 & 常见简单指令走快路径，复杂指令走微序列 & 快路径和微序列是否共享一致的提交边界 \\
\bottomrule
\end{longtable}

\begin{lstlisting}
microinstruction sketch:
  alu_op | src_a | src_b | dst | mem_cmd | pc_cmd | next_micro_pc

LOAD sequence:
  u0: IR  <- MEM[PC],       next = u1
  u1: A   <- R[base],       next = u2
  u2: ADR <- A + imm,       next = u3
  u3: MDR <- MEM[ADR],      next = u4 or wait
  u4: R[dst] <- MDR, PC <- PC+4, next = fetch
\end{lstlisting}

微程序写法有一个重要好处：它迫使作者说明每个周期保存什么中间状态。若某一步需要等待外部存储器，\code{next\_micro\_pc} 可以停留在当前微地址；若访问失败，微序列可以跳到异常入口。这样，等待、错误和普通写回不再混在一句“执行 LOAD”里。

\topic{多周期 CPU 用微操作缩短最长组合路径}单周期 CPU 概念清楚，但周期必须长到容纳取指、译码、ALU、访存和写回的最坏组合路径。多周期 CPU 把同一条指令拆成若干微操作，让每个周期只完成一小段可验收动作。例如 LOAD 可以拆为：PC 取指、译码读寄存器、ALU 形成地址、存储器读、写回寄存器。每一步之间用 IR、MDR、ALUOut 或临时寄存器保存中间结果。

\begin{longtable}{L{0.12\linewidth}L{0.28\linewidth}L{0.3\linewidth}L{0.2\linewidth}}
\toprule
周期 & 微操作 & 打开的硬件路径 & 提交对象 \\
\midrule
T0 & \code{IR = MEM[PC]} & PC 到指令存储器，读响应到 IR & IR \\
T1 & \code{A = R[rs], B = R[rt]} & 指令字段到寄存器堆读地址 & A/B 临时寄存器 \\
T2 & \code{ALUOut = A + imm} & ALU 做地址形成或比较 & ALUOut/Flags \\
T3 & \code{MDR = MEM[ALUOut]} & 地址到数据存储器，读响应到 MDR & MDR \\
T4 & \code{R[rd] = MDR, PC = PC+4} & 写回 mux 和 PC 加法器 & 寄存器、PC \\
\bottomrule
\end{longtable}

这个 trace 说明，多周期设计不是“慢版单周期”，而是把过长的硬件路径切成可计时、可观察、可插入等待的边界。若数据存储器慢，可以只让 T3 等待；若分支比较已经在 T2 得到结果，可以在 T2 或 T3 改写 PC；若发生异常，可以把 PC 选择器改到异常入口，而不是继续提交错误写回。

\topic{先规定一个极小 ISA，程序才能进入 ROM}为了让整机实验真正运行，需要把“指令”压成固定 bit。下面是一种 16-bit 教学编码：高 4 bit 是 opcode，中间字段选择寄存器，低位可解释为寄存器号、立即数或分支偏移。它不是任何真实 ISA 的复制品，但足以把编码、译码、控制字和 ROM 内容连接起来。

\begin{longtable}{L{0.16\linewidth}L{0.22\linewidth}L{0.24\linewidth}L{0.28\linewidth}}
\toprule
指令 & 编码格式 & 语义 & 需要的硬件路径 \\
\midrule
\code{ADD} & \code{0001 d a b ----} & \code{R[d]=R[a]+R[b]} & 双读端口、ALU、写回 \\
\code{LDI} & \code{0010 d imm8} & \code{R[d]=zero_extend(imm8)} & 立即数扩展、写回 \\
\code{LOAD} & \code{0011 d a off4} & \code{R[d]=MEM[R[a]+off]} & 地址形成、读存储、写回 \\
\code{STORE} & \code{0100 s a off4} & \code{MEM[R[a]+off]=R[s]} & 地址形成、写数据、字节使能 \\
\code{JZ} & \code{0101 r off8} & 若 \code{R[r]==0} 则 PC 加偏移 & 比较、符号扩展、PC 选择 \\
\code{JMP} & \code{0110 off12} & PC 加偏移 & 目标形成、PC 写入 \\
\bottomrule
\end{longtable}

这个编码能直接写入 ROM。若 \code{GPIO_IN=0xF004}、\code{GPIO_OUT=0xF000}，程序可以先用 \code{LDI} 装入高地址基址，再循环读取按键、条件跳转、写 LED。教学实现中可把 \code{R0} 固定为 0，也可用一个寄存器保存 I/O 基址。关键不是汇编语法，而是 ROM 中每个 16-bit 字都能被译码成控制信号。

\begin{lstlisting}
tiny ROM program sketch:
  LDI   R1, 0xF0       // R1 holds high part of MMIO base
loop:
  LOAD  R2, [R1 + 4]   // read GPIO_IN
  JZ    R2, off
  LDI   R3, 1
  STORE R3, [R1 + 0]   // write GPIO_OUT = 1
  JMP   loop
off:
  LDI   R3, 0
  STORE R3, [R1 + 0]   // write GPIO_OUT = 0
  JMP   loop
\end{lstlisting}

\topic{从一条语句贯通到硬件事务}为了避免本章变成若干概念并列，可以用一条简单语句贯通整机：\code{if (MMIO[GPIO_IN] != 0) MMIO[GPIO_OUT] = 1;}。在软件层面它像一次条件判断；在硬件层面，它至少包含取指、译码、地址形成、MMIO 读、条件分支、MMIO 写和 PC 更新。每一步都可以写成可验收事务。

\begin{longtable}{L{0.16\linewidth}L{0.38\linewidth}L{0.36\linewidth}}
\toprule
层次 & 发生的动作 & 必须检查的证据 \\
\midrule
取指 & PC 指向 ROM 或 I-cache 中的 LOAD/JZ/STORE 编码 & PC、指令存储片选、IR 内容 \\
译码 & opcode 被译成 \code{mem_read}、\code{pc_sel}、\code{mem_write} 等控制线 & 控制字和寄存器读地址 \\
地址形成 & 基址寄存器与偏移形成 \code{GPIO_IN} 或 \code{GPIO_OUT} 地址 & ALU 输入、ALUOut、地址总线 \\
MMIO 读 & GPIO 控制器返回已同步按键状态 & GPIO 片选、读使能、返回数据稳定 \\
条件判断 & 比较结果决定顺序 PC 或分支目标 & Zero 标志、分支控制、下一 PC \\
MMIO 写 & 输出锁存器接收写数据并驱动 LED & 写使能、字节使能、锁存器输出 \\
\bottomrule
\end{longtable}

这个例子也说明，编译后的指令顺序与外设副作用不能随意重排。普通内存访问可以被 cache、写缓冲和预取优化；但 \code{GPIO_OUT} 写入是外部可见动作，必须按设备属性和顺序规则执行。若系统把 MMIO 当作普通 cacheable 内存，LED 可能不会在预期时刻改变，甚至写入被合并或延迟到错误位置。

\topic{五类基本指令覆盖了 CPU 的主要通路}最小 CPU 不必一开始支持复杂指令，但至少要能解释五类动作：寄存器到寄存器运算、加载、存储、条件分支和无条件跳转。它们共同覆盖寄存器堆、ALU、数据存储器、PC 选择器和标志位。

\begin{longtable}{L{0.14\linewidth}L{0.34\linewidth}L{0.32\linewidth}L{0.1\linewidth}}
\toprule
指令类 & 典型路径 & 必须打开的控制 & 写回对象 \\
\midrule
ADD/SUB & 寄存器读出，经 ALU，结果写回寄存器 & \code{alu_op}、\code{wb_sel=ALU}、\code{reg_write} & 寄存器、Flags \\
LOAD & 基址寄存器加立即数形成地址，存储器读回数据 & \code{alu_op=ADD}、\code{mem_read}、\code{wb_sel=MEM} & 寄存器 \\
STORE & 基址寄存器加立即数形成地址，源寄存器写入存储器 & \code{alu_op=ADD}、\code{mem_write}、\code{store_data_sel} & 存储器或 MMIO \\
BRANCH & 比较源操作数或标志位，条件成立则选择分支目标 & \code{alu_op=CMP}、\code{pc_sel}、\code{pc_write} & PC \\
JUMP & 立即数、寄存器或链接地址决定下一 PC & \code{target_sel}、\code{pc_sel}、可选 \code{reg_write} & PC，可选链接寄存器 \\
\bottomrule
\end{longtable}

以 LOAD 为例，\code{R3 = MEM[R1 + 8]} 不是单纯“读内存”。硬件要先读出 R1，把立即数 8 符号扩展或零扩展到 ALU 输入宽度，ALU 做地址加法，地址进入数据存储路径，读响应回来后经过写回 mux，最后在时钟沿写入 R3。若目标地址落入普通 RAM，读回的是存储阵列内容；若目标地址落入 MMIO 区，读回的是设备寄存器当前状态。地址相同形式，语义由地址译码决定。

\begin{lstlisting}
LOAD R3, [R1 + 8]
  EX:  addr = R1 + sign_extend(8)
  MEM: read target selected by addr
  WB:  R3 = read_data
\end{lstlisting}

STORE 的风险更高，因为它有副作用。 \code{MEM[R1 + 8] = R3} 若落到 RAM，表示保存数据；若落到设备命令寄存器，可能启动设备、清除中断或改变输出引脚。因此 STORE 必须明确字节使能、写数据来源、地址空间属性和完成条件。普通内存写响应表示写事务被接受；设备命令写响应通常不表示设备动作完成。

\begin{lstlisting}
STORE [R1 + 8], R3
  EX:  addr = R1 + sign_extend(8)
  MEM: write_data = R3, byte_en = access_size
       target receives write command
  WB:  no register writeback
\end{lstlisting}

BRANCH 和 JUMP 则说明 PC 是数据通路的一部分。条件分支一般要先得到比较证据，例如 Zero、Sign、Carry 或 Overflow，再由控制器决定 \code{pc_sel}。无条件跳转可以直接选择目标地址；带链接跳转还会把返回地址写入寄存器。此时同一条指令可能同时写 PC 和一个通用寄存器，控制器必须明确两个写边界是否在同一时钟沿提交。

\topic{最小整机实验要先有可观察边界}若把本节落到一块低速实验板，不必立即追求完整 ISA。可以先定义四条动作：\code{ADD}、\code{LOAD}、\code{STORE}、\code{JZ}。指令 ROM 提供编码，4 个寄存器保存通用值，74HC283/74HC86 构成加减核心，SRAM 或寄存器阵列模拟数据存储器，若干 LED 和拨码开关映射成 MMIO。验收重点是每一步都可观察：PC LED 显示当前取指地址，IR LED 或逻辑分析仪显示当前指令，控制线可单步观察，写回和 PC 更新只在时钟沿发生。

\begin{longtable}{L{0.18\linewidth}L{0.38\linewidth}L{0.34\linewidth}}
\toprule
实验对象 & 可用器件或实现 & 验收证据 \\
\midrule
指令 ROM & EEPROM、Flash 或预置拨码/仿真 ROM & 复位后固定地址输出第一条编码 \\
PC/IR & 74HC161、74HC574 或触发器组 & PC 每步按控制更新，IR 保存本周期指令 \\
数据 RAM & SRAM、寄存器阵列或仿真存储器 & LOAD/STORE 地址、数据、写使能可观察 \\
地址译码 & 74HC138、比较器或门阵列 & RAM、LED、按键寄存器不会同时响应 \\
MMIO LED & 锁存器驱动 LED & 写指定地址后 LED 状态在边界改变 \\
MMIO 按键 & 上拉/下拉、施密特、同步器 & 读指定地址返回已同步状态，不直接用生按键 \\
\bottomrule
\end{longtable}

\topic{最小整机的单步调试顺序}搭建教学 CPU 时，最可靠的调试顺序不是一次性运行程序，而是把每个状态边界单独验收。先让复位把 PC 固定到入口地址，再单步确认 ROM 输出第一条指令；随后确认 IR 能锁存该指令，译码能产生正确控制线；再用一条 ADD 验证寄存器读、ALU 和写回；最后才加入 LOAD、STORE、分支和 MMIO。

\begin{longtable}{L{0.16\linewidth}L{0.38\linewidth}L{0.36\linewidth}}
\toprule
调试阶段 & 操作 & 通过标准 \\
\midrule
复位 & 拉低/拉高复位并给一个时钟 & PC 固定为入口，写使能均处于安全值 \\
取指 & 单步一个取指周期 & ROM 片选有效，IR 保存预期编码 \\
ADD & 预置 R1/R2，执行 \code{ADD R3,R1,R2} & R3 在时钟沿变为正确和，PC 前进 \\
LOAD & RAM 某地址预置数据，执行加载 & 地址、读使能、写回来源均正确 \\
STORE & 执行写 RAM 或写 LED 寄存器 & 只有目标片选有效，非目标不被改写 \\
分支 & 改变 Zero 条件重复执行 & PC 在顺序地址和分支目标之间正确选择 \\
\bottomrule
\end{longtable}

这种验收方式也适用于 FPGA。逻辑分析仪、仿真波形或片上 ILA 应至少抓取 PC、IR、控制字、ALU 输入输出、内存地址、读写使能、写回数据和寄存器写地址。若只看最终 LED 是否闪烁，错误可能被循环掩盖；若能看到每条指令的提交边界，错误会落到具体周期。

\topic{流水线把吞吐率问题提前暴露出来}一旦把取指、译码、执行、访存和写回拆成流水级，CPU 的平均吞吐率可以提高，但控制复杂度也立刻上升。相邻指令可能争用同一个存储器端口，后一条指令可能读取前一条尚未写回的寄存器，分支可能让已经取到的指令变成错误路径。这些问题统称结构冒险、数据冒险和控制冒险。

\begin{longtable}{L{0.18\linewidth}L{0.34\linewidth}L{0.38\linewidth}}
\toprule
冒险类型 & 例子 & 硬件处理 \\
\midrule
结构冒险 & 取指和 LOAD 同时争用单端口存储器 & 分离指令/数据存储，或插入停顿 \\
数据冒险 & \code{ADD R1,...} 后立刻 \code{SUB R2,R1,...} & 旁路/转发，或等待写回后再读 \\
控制冒险 & 条件分支结果未出时已取后续指令 & 预测、延迟槽、停顿或冲刷流水线 \\
异常冒险 & 较晚指令先完成，较早指令发生异常 & 提交队列或按序提交保证精确状态 \\
\bottomrule
\end{longtable}

因此，现代芯片低时延不只是“流水线更深”。流水线必须配套转发网络、停顿控制、冲刷控制和精确异常边界。否则，吞吐率提高会以错误状态提交为代价。学习最小 CPU 时提前看到这些冒险，可以防止把流水线误解成简单地把框图切成几段。

\topic{本章到此只验收 CPU 内部 trace}最小 CPU 章节到这里应停在内部状态边界：PC、IR、寄存器堆、ALU、写回 mux、控制字、等待状态和流水线冒险。只要读者能把每条指令写成“读旧状态、经过组合路径、在时钟沿提交新状态”的 trace，本章目标就已经完成。

下一章再把这个 CPU 接到主存、总线、MMIO、DMA 和中断，讨论地址译码、等待、错误、cache 可见性和设备副作用。第六章再用 6502、Z80、8051、8086、80386、80486/Pentium 和现代 SoC 观察真实芯片如何把这些边界逐步扩展到整机平台。这样分章后，CPU 内部 trace、整机事务、经典芯片三类问题不会混在同一节里。



\section*{章末练习与验收任务}

本章练习不要求先追求完整 ISA，而要求把每条指令写成可检查的硬件 trace。答案必须同时说明控制线、数据路径和提交边界。

\begin{longtable}{L{0.2\linewidth}L{0.48\linewidth}L{0.22\linewidth}}
\toprule
任务 & 要完成的产物 & 验收标准 \\
\midrule
PC 取指 & 写出 PC、指令存储器、IR 和 \code{pc_write/ir_write} 的周期关系 & 能说明哪一拍保存指令 \\
ADD trace & 为 \code{R3=R1+R2} 列出读地址、ALU 操作、写回选择和写使能 & 写回只在目标时钟沿发生 \\
LOAD trace & 把 LOAD 拆成地址形成、存储读、MDR 保存、寄存器写回 & 能处理存储器未 ready \\
STORE trace & 说明地址、写数据、字节使能、写控制和无寄存器写回 & 能区分 RAM 写和 MMIO 副作用 \\
JZ trace & 说明 Zero 标志如何影响下一 PC & 分支目标和顺序 PC 选择明确 \\
控制字 & 设计一个包含 \code{alu_op/wb_sel/reg_write/mem_read/mem_write/pc_sel} 的控制字 & 每一位能追到实际电路 \\
微程序 & 为 LOAD 写出 5 个微步骤和等待状态处理 & 慢存储器只拖住访存阶段 \\
流水线冒险 & 说明结构、数据、控制三类冒险各破坏什么边界 & 能提出停顿、转发或冲刷策略 \\
\bottomrule
\end{longtable}

\section*{参考解答与验收要点}

\begin{longtable}{L{0.18\linewidth}L{0.5\linewidth}L{0.22\linewidth}}
\toprule
任务 & 参考解答 & 必须保留的验收点 \\
\midrule
PC 取指 & PC 输出取指地址，指令存储器返回编码，\code{ir_write} 在取指边界保存编码；下一 PC 可由 \code{pc+size}、分支目标或异常入口选择。 & PC 和 IR 都是状态元件。 \\
ADD trace & 读端口选择 R1/R2，ALU 选择 ADD，写回 mux 选择 ALUOut，写地址为 R3，\code{reg_write=1} 只在写回边界有效。 & 算对不等于写对。 \\
LOAD trace & EX 阶段形成地址，MEM 阶段发起读，目标未 ready 时保持地址和控制，响应后保存 MDR，WB 阶段写回目标寄存器。 & 等待不能污染其他状态。 \\
STORE trace & STORE 使用 ALU 形成地址，源寄存器提供写数据，打开字节使能和写控制，不打开寄存器写回。若地址命中 MMIO，副作用由设备规格决定。 & STORE 是有副作用的事务。 \\
JZ trace & 比较或 ALU 结果产生 Zero，控制器根据 Zero 选择顺序 PC 或分支目标；PC 只在提交边界更新。 & 条件输入必须在采样前稳定。 \\
控制字 & 控制字字段应覆盖 PC、IR、寄存器堆、ALU、存储器和下一状态逻辑；任何字段没有连接对象都不是硬件控制。 & 控制字是电线集合，不是注释。 \\
微程序 & LOAD 可写为取指、读寄存器、地址形成、访存等待、写回；访存未完成时微地址保持在等待状态，错误时跳异常入口。 & 微步骤提供可观察边界。 \\
流水线冒险 & 结构冒险来自资源争用，数据冒险来自结果未写回，控制冒险来自下一 PC 尚未确定。停顿保护边界，转发减少等待，冲刷丢弃错误路径。 & 提高吞吐不能破坏提交顺序。 \\
\bottomrule
\end{longtable}

\begin{classicnote}
最小 CPU 的经典读法是逐周期 trace。不要只画“PC 连到内存、ALU 连到寄存器”的框图；要写清楚每条指令在哪个周期读旧状态、在哪条组合路径上传播、在哪个时钟沿提交新状态。后续主存、总线、I/O 和经典芯片只是在这个 trace 外面继续增加事务边界。
\end{classicnote}
